% Source-derived chapter 00: Abstract and front matter
% Original source: riemann-roch-polynomials-hyperkahler-fourfolds
\chapter{Abstract and front matter}
\label{riemann-roch-polynomials-hyperkahler-fourfolds-chapter-00}
\source{riemann-roch-polynomials-hyperkahler-fourfolds}

\begin{remark}[Source section]
\label{riemann-roch-polynomials-hyperkahler-fourfolds-chapter-00-source}
\source{riemann-roch-polynomials-hyperkahler-fourfolds}
This chapter records the source section; no Lean declaration has yet been checked for this material.
\notready
\end{remark}

% The source section title was: Abstract and front matter





















\begin{abstract}
We prove that a hyper-K\"ahler fourfold satisfying a mild topological assumption is of K3$^{[2]}$ deformation type.\
This proves in particular a conjecture of O'Grady stating that hyper-K\"ahler fourfolds of K3$^{[2]}$ numerical type are of K3$^{[2]}$ deformation type.\
Our topological assumption concerns the existence of two integral degree-2 cohomology classes satisfying certain numerical intersection conditions.

There are two main ingredients in the proof.\
We first prove a topological version of the statement, by showing that our topological assumption forces the Betti numbers, the Fujiki constant, and the Huybrechts--Riemann--Roch polynomial of the hyper-K\"ahler fourfold to be the same as those of K3$^{[2]}$ hyper-K\"ahler fourfolds.\
The key part of the article is then to prove the hyper-K\"ahler SYZ conjecture for hyper-K\"ahler fourfolds for divisor classes satisfying the numerical condition mentioned above.
\end{abstract}
